% =============================================================================
%  RL-H5-01 | Figur uji sintetis
% -----------------------------------------------------------------------------
%  Role    : RL — Research Lead (kurasi; MLE menyiapkan data)
%  Sprint  : H5 | PARALEL
%  Ticket  : (isi dari papan kerja)          Status: STARTER ONLY
%  Target  : reports/figures/synth_lag_recovery.pdf,
%            reports/figures/synth_edge_heatmap.pdf,
%            reports/figures/captions.tex
%  Input   : reports/synth/recovery_metrics.csv, lag_true_vs_hat.parquet,
%            A_by_regime.npz
%  Output  : figur + caption
%  Handoff : laporan
%  DoD     : [ ] jeda benar vs terpelajar per edge, dipisah regime
%            [ ] heatmap A benar vs terpelajar (34x34) untuk 2 regime
%            [ ] caption menyebut ambang protokol dan lolos/tidak
% =============================================================================
\documentclass[11pt,a4paper]{article}
% =============================================================================
%  00_preamble.tex | shared preamble untuk semua deliverable RL
%  Job     : bukan job card; dipakai oleh RL-H1-01 .. RL-H7-01
%  Target  : docs/ (salin bersama file .tex)
%  Catatan : kompilasi dengan pdflatex. Untuk XeLaTeX/LuaLaTeX, hapus inputenc.
% =============================================================================

\usepackage[utf8]{inputenc}
\usepackage[T1]{fontenc}
\usepackage{lmodern}
\usepackage[a4paper,margin=2.4cm]{geometry}
\usepackage{amsmath,amssymb}
\usepackage{booktabs}
\usepackage{longtable}
\usepackage{array}
\usepackage{graphicx}
\usepackage{xcolor}
\usepackage{enumitem}
\usepackage{float}
\usepackage{caption}
\usepackage{url}
\usepackage{microtype}
\usepackage[colorlinks=true,linkcolor=blue!50!black,urlcolor=blue!50!black,citecolor=blue!50!black]{hyperref}

% --- helper wajib -----------------------------------------------------------
% Path/nama file aman (underscore, garis miring) di dalam teks.
\newcommand{\fpath}[1]{\texttt{\detokenize{#1}}}
% TODO yang mencolok. Jangan dihapus sebelum job selesai.
\newcommand{\todo}[1]{\textcolor{red!75!black}{\textbf{[TODO: #1]}}}
\newcommand{\jobid}[1]{\textsf{\textbf{#1}}}
\newcommand{\runid}[1]{\texttt{\detokenize{#1}}}
\newcommand{\claim}[1]{\textsf{#1}}
\newcommand{\good}{\textcolor{green!45!black}{$\checkmark$}}
\newcommand{\bad}{\textcolor{red!75!black}{$\times$}}
\newcommand{\maybe}{\textcolor{orange!80!black}{$\sim$}}

\newcolumntype{L}[1]{>{\raggedright\arraybackslash}p{#1}}
\newcolumntype{C}[1]{>{\centering\arraybackslash}p{#1}}

\setlength{\parindent}{0pt}
\setlength{\parskip}{4pt}
\captionsetup{font=small,labelfont=bf}
\setlist[itemize]{leftmargin=1.2em,itemsep=2pt,topsep=2pt}
\setlist[enumerate]{leftmargin=1.5em,itemsep=2pt,topsep=2pt}

% Judul job yang konsisten di semua dokumen.
\newcommand{\jobheader}[4]{%
  \begin{center}
  {\Large\bfseries #1}\\[2pt]
  {\large #2}\\[4pt]
  \begin{tabular}{@{}ll@{}}
  \toprule
  Role & #3 \\
  Sprint & #4 \\
  \bottomrule
  \end{tabular}
  \end{center}
}


\title{\jobid{RL-H5-01} | Figur Uji Sintetis\\
\large Apakah Modul 1 menemukan kembali pergeseran jeda?}
\author{RL \\ \small RENARO Lab}
\date{\todo{tanggal H5} \quad | \quad \textcolor{red!75!black}{STARTER}}

\begin{document}
\maketitle

\section{Ambang protokol}
\begin{table}[H]
\centering
\small
\begin{tabular}{@{}L{6.0cm}C{2.4cm}C{2.4cm}L{3.4cm}@{}}
\toprule
Metrik & Ambang & Nilai & Lolos? \\
\midrule
Akurasi argmax-jeda (edge benar) & $\ge 0{,}80$ & \todo{} & \todo{} \\
Spearman massa-jeda vs jeda benar & $\ge 0{,}70$ & \todo{} & \todo{} \\
$\Delta$lag (kering $-$ normal) & mendekati 2 & \todo{} & \todo{} \\
Kontrol tanpa iklim & $\Delta$lag $\approx 0$ & \todo{} & \todo{} \\
Kontrol \texttt{unconditioned} & $\Delta$lag $\approx 0$ & \todo{} & \todo{} \\
\bottomrule
\end{tabular}
\end{table}

\section{Figur 1: recovery jeda}
\begin{figure}[H]
\centering
\fbox{\parbox[c][7.2cm][c]{0.88\textwidth}{\centering
\color{gray} PLACEHOLDER \\[4pt]
\fpath{reports/figures/synth_lag_recovery.pdf} \\[4pt]
(TODO: jeda benar vs jeda terpelajar per edge; titik/hexbin; \\ warna = regime normal vs kering; garis $y=x$)}}
\caption{\todo{Caption wajib: data variant, model variant, 5 seed, ambang protokol,
dan kesimpulan lolos/tidak. Sebut bahwa sumbu-x adalah jeda yang ditanam (pentad)
dan sumbu-y adalah argmax massa perhatian.}}
\end{figure}

\section{Figur 2: heatmap A benar vs terpelajar}
\begin{figure}[H]
\centering
\fbox{\parbox[c][7.2cm][c]{0.88\textwidth}{\centering
\color{gray} PLACEHOLDER \\[4pt]
\fpath{reports/figures/synth_edge_heatmap.pdf} \\[4pt]
(TODO: 2x2 panel: A benar, A terpelajar normal, A benar, A terpelajar kering; \\ skala warna sama; diagonal dikosongkan)}}
\caption{\todo{Caption wajib: 34 node, arah $A[i,j]=j\to i$, rata-rata 5 seed,
dan catatan bahwa ini dependensi prediktif, bukan arus pasokan.}}
\end{figure}

\section{Catatan interpretasi}
\todo{Tulis 3--5 kalimat: apakah pergeseran ditemukan, di regime mana gagal,
apakah kontrol bersih, dan apa implikasinya untuk klaim metode. Jangan
menggeneralisasi dari satu seed.}

\end{document}
